\documentclass[10pt,a4paper]{amsart}
\usepackage[margin=19mm]{geometry}
\usepackage{amsmath,amssymb,mathtools}
\usepackage[scr=rsfs,cal=euler]{mathalfa}
\usepackage{xfrac}
\usepackage[unicode,hidelinks,bookmarks=false]{hyperref}
\newcommand{\ie}{i.e.\ }
\newcommand{\cf}{cf.\ }
\newcommand{\resp}{respectively\ }
\newcommand{\Subsection}{\subsection}
\providecommand{\nobreakdash}{\nobreak\hbox{-}}
\setlength{\emergencystretch}{2em}
\multlinegap0pt
\usepackage{tikz}
\usepackage{enumerate}
\usepackage{amssymb}
\newtheorem{thm}{Theorem}
\newcommand{\N}{\mathbf{n}}
\title{Higher-rank graphs and the graded $K$-theory of Kumjian--Pask Algebras}
\author{Roozbeh Hazrat, Promit Mukherjee, David Pask, Sujit Kumar Sardar}
\date{Source: arXiv:2507.19879v3 (CC BY 4.0)}
\begin{document}
\maketitle

\noindent\emph{Source and licence.} Higher-rank graphs and the graded $K$-theory of Kumjian--Pask Algebras. Version arXiv:2507.19879v3 (CC BY 4.0), distributed by arXiv under the Creative Commons Attribution 4.0 International licence, http://creativecommons.org/licenses/by/4.0/, as recorded in the arXiv licence metadata for this exact version and shown on the abstract page https://arxiv.org/abs/2507.19879v3. Full licence notice: \texttt{licenses/arxiv-2507.19879-CC-BY-4.0.txt}.\par\smallskip

\noindent\textit{Editorial scope.} Counted unit: the article's thm monoid (source0.tex lines 629--639). The article's complete original proof of it is delivered with it (source0.tex lines 640--717, one \texttt{proof} environment of 78 lines). No mathematics is written, repaired or re-derived: the statement and the proof are the article's own lines, in the article's own notation, and no retained line is substituted or re-flowed.  No further source-local context is needed: every cross-reference of every retained line resolves inside the two delivered spans.  Nothing was omitted from any retained span, so the package carries no omission record.  No retained line contains a citation, so the wrapper carries no bibliography and no bibliography entry.  No retained line contains a figure environment, an image-inclusion command or drawing code, so no image is delivered and the package declares no asset dependency.  Editorial formatting only, in addition: an AMS \texttt{amsart} preamble that restates the article's own declarations and macros used by the retained lines, namely the environments \texttt{thm} on separate counters, together with the standard packages those lines need.  The retained environments print in this document as Theorem 1 in order of appearance, whereas the article prints the same items with the numbering of its own full document.  The complete native source of the article as distributed in the arXiv e-print for this exact version is bundled unchanged as \texttt{sources/182246/source0.tex}.
\begin{thm}\label{th in-splitting preserves talented monoid}
Let $\Lambda$ be a row-finite $k$-graph with no sources and $\Lambda_I$ the in-split of $\Lambda$ at a vertex $v$ with respect to the partition $v\Lambda^\mathbf{1}=\mathcal{E}_1\cup \mathcal{E}_2$. The map $\phi_I: T_\Lambda \longrightarrow T_{\Lambda_I}$ defined on generators as
$$\phi_I(w(\N)):=
	\left\{
	\begin{array}{ll}
		w(\N)  & \mbox{if } w\neq v, \\
		v^1(\N)+v^2(\N)  & \mbox{if } w=v
	\end{array}
	\right.$$ 
for all $w\in \Lambda^0$ and $\N\in \mathbb{Z}^k$, is a $\mathbb{Z}^k$-monoid isomorphism. 
\end{thm}

\begin{proof}
We first show that $\phi_I$ is well-defined. For this, we need to show that $\phi_I$ preserves the relations \[u(\N)=\displaystyle{\sum_{\lambda\in u\Lambda^{\mathbf{e}_i}}} s(\lambda)(\N+\mathbf{e}_i),\] for all $u\in \Lambda^0$, $\N\in \mathbb{Z}^k$ and $i=1,2,\ldots,k$. Note that \[\phi_I\left(\displaystyle{\sum_{\lambda\in u\Lambda^{\mathbf{e}_i}}} s(\lambda)(\N+\mathbf{e}_i)\right)=\displaystyle{\sum_{\lambda\in u\Lambda^{\mathbf{e}_i}}}\phi_I(s(\lambda)(\N+\mathbf{e}_i))=\displaystyle{\sum_{\substack{\lambda\in u\Lambda^{\mathbf{e}_i}\\s(\lambda)\neq v}}} s(\lambda)(\N+\mathbf{e}_i)+\displaystyle{\sum_{\lambda\in u\Lambda^{\mathbf{e}_i}v}} (v^1(\N+\mathbf{e}_i)+v^2(\N+\mathbf{e}_i)).\] If $u\neq v$, then it implies 
{\allowdisplaybreaks
\begin{align*}
&\phi_I\left(\displaystyle{\sum_{\lambda\in u\Lambda^{\mathbf{e}_i}}} s(\lambda)(\N+\mathbf{e}_i)\right)\\
=& \displaystyle{\sum_{\substack{\alpha\in u\Lambda_I^{\mathbf{e}_i}\\s(par(\alpha))\neq v}}} s_I(\alpha)(\N+\mathbf{e}_i)+\displaystyle{\sum_{\beta\in u\Lambda_I^{\mathbf{e}_i}v^1}} s_I(\beta)(\N+\mathbf{e}_i)+\displaystyle{\sum_{\gamma\in u\Lambda_I^{\mathbf{e}_i}v^2}} s_I(\gamma)(\N+\mathbf{e}_i)\\
=& \displaystyle{\sum_{h\in u\Lambda_I^{\mathbf{e}_i}}} s_I(h)(\N+\mathbf{e}_i)\\
=&~ u(\N)=\phi_I(u(\N)),
\end{align*}
}
whereas 
\noindent
{\allowdisplaybreaks
\begin{align*}
&\phi_I\left(\displaystyle{\sum_{\lambda\in v\Lambda^{\mathbf{e}_i}}} s(\lambda)(\N+\mathbf{e}_i)\right)\\
=& \displaystyle{\sum_{\substack{\lambda\in v\Lambda^{\mathbf{e}_i}\cap \mathcal{E}_1\\s(\lambda)\neq v}}} s(\lambda)(\N+\mathbf{e}_i)+\displaystyle{\sum_{\substack{\lambda\in v\Lambda^{\mathbf{e}_i}\cap \mathcal{E}_2\\s(\lambda)\neq v}}} s(\lambda)(\N+\mathbf{e}_i)+\displaystyle{\sum_{\lambda\in v\Lambda^{\mathbf{e}_i}v\cap \mathcal{E}_1}}(v^1(\N+\mathbf{e}_i)+v^2(\N+\mathbf{e}_i))\\
&+\displaystyle{\sum_{\lambda\in v\Lambda^{\mathbf{e}_i}v\cap \mathcal{E}_2}} (v^1(\N+\mathbf{e}_i)+v^2(\N+\mathbf{e}_i))\\
=& \left(\displaystyle{\sum_{\substack{\lambda\in v\Lambda^{\mathbf{e}_i}\cap \mathcal{E}_1\\s(\lambda)\neq v}}} s(\lambda)(\N+\mathbf{e}_i)+\displaystyle{\sum_{\lambda\in v\Lambda^{\mathbf{e}_i}v \cap \mathcal{E}_1}}v^1(\N+\mathbf{e}_i)+\displaystyle{\sum_{\lambda\in v\Lambda^{\mathbf{e}_i}v \cap \mathcal{E}_1}}v^2(\N+\mathbf{e}_i)\right)\\
&+\left(\displaystyle{\sum_{\substack{\lambda\in v\Lambda^{\mathbf{e}_i}\cap \mathcal{E}_2\\s(\lambda)\neq v}}} s(\lambda)(\N+\mathbf{e}_i)+\displaystyle{\sum_{\lambda\in v\Lambda^{\mathbf{e}_i}v \cap \mathcal{E}_2}}v^1(\N+\mathbf{e}_i)+\displaystyle{\sum_{\lambda\in v\Lambda^{\mathbf{e}_i}v \cap \mathcal{E}_2}}v^2(\N+\mathbf{e}_i)\right)\\
=& \left(\displaystyle{\sum_{\substack{\alpha\in v^1\Lambda_I^{\mathbf{e}_i}\\s(par(\alpha))\neq v}}} s_I(\alpha)(\N+\mathbf{e}_i)+\displaystyle{\sum_{\beta\in v^1\Lambda_I^{\mathbf{e}_i}v^1}} s_I(\beta)(\N+\mathbf{e}_i)+\displaystyle{\sum_{\gamma\in v^1\Lambda_I^{\mathbf{e}_i}v^2}} s_I(\gamma)(\N+\mathbf{e}_i)\right)\\
&+\left(\displaystyle{\sum_{\substack{\alpha'\in v^2\Lambda_I^{\mathbf{e}_i}\\s(par(\alpha'))\neq v}}} s_I(\alpha')(\N+\mathbf{e}_i)+\displaystyle{\sum_{\beta'\in v^2\Lambda_I^{\mathbf{e}_i}v^1}} s_I(\beta')(\N+\mathbf{e}_i)+\displaystyle{\sum_{\gamma'\in v^2\Lambda_I^{\mathbf{e}_i}v^2}} s_I(\gamma')(\N+\mathbf{e}_i)\right)\\
=& \displaystyle{\sum_{g\in v^1\Lambda_I^{\mathbf{e}_i}}} s_I(g)(\N+\mathbf{e}_i)+\displaystyle{\sum_{h\in v^2\Lambda_I^{\mathbf{e}_i}}} s_I(h)(\N+\mathbf{e}_i)\\
=&~ v^1(\N)+v^2(\N)=\phi_I(v(\N)).
\end{align*}
} 
Thus, $\phi_I$ is a well-defined monoid homomorphism. We now define an inverse for $\phi_I$. For this, we fix $j\in \{1,2,\ldots,k\}$ and define a map $\psi:T_{\Lambda_I}\longrightarrow T_\Lambda$ on generators by 
$$\psi(u(\N)):=
	\left\{
	\begin{array}{ll}
		\hspace{1.25cm} u(\N)  & \mbox{if } u\neq v^1,v^2, \\
		\displaystyle{\sum_{f\in \mathcal{E}_i\cap \Lambda^{\mathbf{e}_j}}} s(f)(\N+\mathbf{e}_j)  & \mbox{if } u=v^i~\mbox{for}~i=1,2.
	\end{array}
	\right.$$ 
Again we first need to check that $\psi$ is well-defined. If $u\neq v^1,v^2$, then for each $i=1,2,\ldots,k$ and $\N\in \mathbb{Z}^k$, we have 
{\allowdisplaybreaks
\begin{align*}
&\psi\left(\displaystyle{\sum_{\alpha\in u\Lambda_I^{\mathbf{e}_i}}} s_I(\alpha)(\N+\mathbf{e}_i)\right)\\
=& \displaystyle{\sum_{\alpha\in u\Lambda_I^{\mathbf{e}_i}}} \psi(s_I(\alpha)(\N+\mathbf{e}_i))\\
=& \displaystyle{\sum_{\substack{\alpha\in u\Lambda_I^{\mathbf{e}_i}\\s_I(\alpha)\neq v^1,v^2}}} s_I(\alpha)(\N+\mathbf{e}_i)+\displaystyle{\sum_{\alpha\in u\Lambda_I^{\mathbf{e}_i}v^1}} \left(\displaystyle{\sum_{f\in \mathcal{E}_1\cap \Lambda^{\mathbf{e}_j}}} s(f)(\N+\mathbf{e}_i+\mathbf{e}_j)\right)+\displaystyle{\sum_{\alpha\in u\Lambda_I^{\mathbf{e}_i}v^2}} \left(\displaystyle{\sum_{f\in \mathcal{E}_2\cap \Lambda^{\mathbf{e}_j}}} s(f)(\N+\mathbf{e}_i+\mathbf{e}_j)\right)\\
=& \displaystyle{\sum_{\substack{\alpha\in u\Lambda_I^{\mathbf{e}_i}\\par(\alpha)\neq v}}}s(par(\alpha))(\N+\mathbf{e}_i)+\displaystyle{\sum_{\lambda\in u\Lambda^{\mathbf{e}_i}v}}\left(\displaystyle{\sum_{f\in \mathcal{E}_1\cap \Lambda^{\mathbf{e}_j}}} s(f)(\N+\mathbf{e}_i+\mathbf{e}_j)+\displaystyle{\sum_{f\in \mathcal{E}_2\cap \Lambda^{\mathbf{e}_j}}} s(f)(\N+\mathbf{e}_i+\mathbf{e}_j)\right)\\
=& \displaystyle{\sum_{\substack{\lambda\in u\Lambda^{\mathbf{e}_i}\\s(\lambda)\neq v}}} s(\lambda)(\N+\mathbf{e}_i)+\displaystyle{\sum_{\lambda\in u\Lambda^{\mathbf{e}_i}v}}\left(\displaystyle{\sum_{f\in v\Lambda^{\mathbf{e}_j}}} s(f)(\N+\mathbf{e}_i+\mathbf{e}_j)\right)\\
=& \displaystyle{\sum_{\substack{\lambda\in u\Lambda^{\mathbf{e}_i}\\s(\lambda)\neq v}}}s(\lambda)(\N+\mathbf{e}_i)+\displaystyle{\sum_{\lambda\in u\Lambda^{\mathbf{e}_i}v}} v(\N+\mathbf{e}_i) ~~(\text{using relations of}~ T_\Lambda)\\
=& \displaystyle{\sum_{\lambda\in u\Lambda^{\mathbf{e}_i}}} s(\lambda)(\N+\mathbf{e}_i)=u(\N)=\psi(u(\N)).
\end{align*}
For $t=1,2$ 
\begin{align*}
&\psi\left(\displaystyle{\sum_{\alpha\in v^t\Lambda_I^{\mathbf{e}_i}}} s_I(\alpha)(\N+\mathbf{e}_i)\right)\\
=&  \displaystyle{\sum_{\substack{\alpha\in v^t\Lambda_I^{\mathbf{e}_i}\\s_I(\alpha)\neq v^1,v^2}}} s_I(\alpha)(\N+\mathbf{e}_i)+\displaystyle{\sum_{\alpha\in v^t\Lambda_I^{\mathbf{e}_i}v^1}} \left(\displaystyle{\sum_{f\in \mathcal{E}_1\cap \Lambda^{\mathbf{e}_j}}} s(f)(\N+\mathbf{e}_i+\mathbf{e}_j)\right)+\displaystyle{\sum_{\alpha\in v^t\Lambda_I^{\mathbf{e}_i}v^2}} \left(\displaystyle{\sum_{f\in \mathcal{E}_2\cap \Lambda^{\mathbf{e}_j}}} s(f)(\N+\mathbf{e}_i+\mathbf{e}_j)\right)\\
=& \displaystyle{\sum_{\substack{\alpha\in v^t\Lambda_I^{\mathbf{e}_i}\\par(\alpha)\neq v}}}s(par(\alpha))(\N+\mathbf{e}_i)+\displaystyle{\sum_{\lambda\in v\Lambda^{\mathbf{e}_i}v\cap \mathcal{E}_t}}\left(\displaystyle{\sum_{f\in \mathcal{E}_1\cap \Lambda^{\mathbf{e}_j}}} s(f)(\N+\mathbf{e}_i+\mathbf{e}_j)+\displaystyle{\sum_{f\in \mathcal{E}_2\cap \Lambda^{\mathbf{e}_j}}} s(f)(\N+\mathbf{e}_i+\mathbf{e}_j)\right)\\
=& \displaystyle{\sum_{\substack{\lambda\in v\Lambda^{\mathbf{e}_i}\cap \mathcal{E}_t\\s(\lambda)\neq v}}} s(\lambda)(\N+\mathbf{e}_i)+\displaystyle{\sum_{\lambda\in v\Lambda^{\mathbf{e}_i}v\cap \mathcal{E}_t}}\left(\displaystyle{\sum_{f\in v\Lambda^{\mathbf{e}_j}}} s(f)(\N+\mathbf{e}_i+\mathbf{e}_j)\right)\\
=& \displaystyle{\sum_{\substack{\lambda\in v\Lambda^{\mathbf{e}_i}\cap \mathcal{E}_t\\s(\lambda)\neq v}}}s(\lambda)(\N+\mathbf{e}_i)+\displaystyle{\sum_{\lambda\in v\Lambda^{\mathbf{e}_i}v\cap \mathcal{E}_t}} v(\N+\mathbf{e}_i) ~~(\text{using relations of}~ T_\Lambda)\\
=& \displaystyle{\sum_{\lambda\in \mathcal{E}_t \cap \Lambda^{\mathbf{e}_i}}} s(\lambda)(\N+\mathbf{e}_i).
\end{align*}
}
The second equality above follows since the parent of any $\alpha\in v^t\Lambda_I^{\mathbf{e}_i}v^1$ (or $v^t\Lambda_I^{\mathbf{e}_i}v^2$) is a loop at $v$ of degree $\mathbf{e}_i$, contained in the partition set $\mathcal{E}_t$; and hence both $v^t\Lambda_I^{\mathbf{e}_i}v^1$ and $v^t\Lambda_I^{\mathbf{e}_i}v^2$ are in bijection with $v\Lambda^{\mathbf{e}_i}v \cap \mathcal{E}_t$. The crucial part which remains to prove in order to complete the verification is that the equality
\begin{equation}\label{welldefined}
\displaystyle{\sum_{\lambda\in \mathcal{E}_t \cap \Lambda^{\mathbf{e}_i}}} s(\lambda)(\N+\mathbf{e}_i)=\displaystyle{\sum_{\mu\in \mathcal{E}_t \cap \Lambda^{\mathbf{e}_j}}} s(\mu)(\N+\mathbf{e}_j)    
\end{equation}
holds in $T_\Lambda$ for each $i=1,2,\ldots,k$. There is nothing to show if $i=j$. So assume that $i\neq j$. Then, in $T_\Lambda$,
\begin{align*}
\displaystyle{\sum_{\lambda\in \mathcal{E}_t\cap \Lambda^{\mathbf{e}_i}}} s(\lambda)(\N+\mathbf{e}_i)&=\displaystyle{\sum_{\lambda\in \mathcal{E}_t\cap \Lambda^{\mathbf{e}_i}}} \left(\displaystyle{\sum_{\alpha\in s(\lambda)\Lambda^{\mathbf{e}_j}}} s(\alpha)(\N+\mathbf{e}_i+\mathbf{e}_j)\right)\\
&= \displaystyle{\sum_{\substack{\gamma\in v\Lambda^{\mathbf{e}_i+\mathbf{e}_j}\\\gamma(0,\mathbf{e}_i)\in \mathcal{E}_t}}} s(\gamma)(\N+\mathbf{e}_i+\mathbf{e}_j)\\
&= \displaystyle{\sum_{\substack{\gamma\in v\Lambda^{\mathbf{e}_i+\mathbf{e}_j}\\\gamma(0,\mathbf{e}_j)\in \mathcal{E}_t}}} s(\gamma)(\N+\mathbf{e}_i+\mathbf{e}_j)\\
&= \displaystyle{\sum_{\mu\in \mathcal{E}_t\cap \Lambda^{\mathbf{e}_j}}} \left(\displaystyle{\sum_{\beta\in s(\mu)\Lambda^{\mathbf{e}_i}}} s(\beta)(\N+\mathbf{e}_j+\mathbf{e}_i)\right)=\displaystyle{\sum_{\mu\in \mathcal{E}_t\cap \Lambda^{\mathbf{e}_j}}} s(\mu)(\N+\mathbf{e}_j).
\end{align*}
The second and fourth equalities are self-explanatory. For the third one, note that whenever $\gamma=\lambda\alpha\in v\Lambda^{\mathbf{e}_i+\mathbf{e}_j}$ and $\lambda\in \mathcal{E}_t\cap \Lambda^{\mathbf{e}_i}$, by the factorization property of $\Lambda$, there exist unique $\mu\in v\Lambda^{\mathbf{e}_j}$ and $\beta\in s(\mu)\Lambda^{\mathbf{e}_i}$ such that $\gamma=\mu\beta$. Then, by the pairing condition, $\lambda$ and $\mu$ are in the same partition set, whence $\mu\in \mathcal{E}_t$. Thus $\gamma(0,\mathbf{e}_i)\in \mathcal{E}_t$ if and only if $\gamma(0,\mathbf{e}_j)\in \mathcal{E}_t$. 

Now using $(\ref{welldefined})$, we have \[\psi\left(\displaystyle{\sum_{\alpha\in v^t\Lambda_I^{\mathbf{e}_i}}} s_I(\alpha)(\N+\mathbf{e}_i)\right)=\displaystyle{\sum_{\lambda\in \mathcal{E}_t\cap \Lambda^{\mathbf{e}_i}}} s(\lambda)(\N+\mathbf{e}_i)=\displaystyle{\sum_{\mu\in \mathcal{E}_t\cap \Lambda^{\mathbf{e}_j}}} s(\mu)(\N+\mathbf{e}_j)=\psi(v^t(\N)),\] showing that $\psi$ is well-defined. To complete the proof, we observe that $\phi_I$ and $\psi$ are mutually inverse homomorphisms. If $w\neq v$ then $\psi(\phi_I(w(\N)))=\psi(w(\N))=w(\N)$. Again, \[\psi(\phi_I(v(\N)))=\psi(v^1(\N)+v^2(\N))=\displaystyle{\sum_{f\in \mathcal{E}_1\cap \Lambda^{\mathbf{e}_j}}} s(f)(\N+\mathbf{e}_j)+\displaystyle{\sum_{g\in \mathcal{E}_2\cap \Lambda^{\mathbf{e}_j}}} s(g)(\N+\mathbf{e}_j)=\displaystyle{\sum_{h\in v\Lambda^{\mathbf{e}_j}}} s(h)(\N+\mathbf{e}_j)=v(\N).\] On the other hand, for $u\neq v^1,v^2$, $\phi_I(\psi(u(\N)))=\phi_I(u(\N))=u(\N)$ and for $t=1,2$,
{\allowdisplaybreaks
\begin{align*}
\phi_I(\psi(v^t(\N)))&=\phi_I\left(\displaystyle{\sum_{f\in \mathcal{E}_t\cap \Lambda^{\mathbf{e}_j}}} s(f)(\N+\mathbf{e}_j)\right)\\
&= \displaystyle{\sum_{\substack{f\in \mathcal{E}_t\cap \Lambda^{\mathbf{e}_j}\\s(f)\neq v}}} s(f)(\N+\mathbf{e}_j)+\displaystyle{\sum_{f\in \mathcal{E}_t\cap \Lambda^{\mathbf{e}_j}v}} (v^1(\N+\mathbf{e}_j)+v^2(\N+\mathbf{e}_j))\\
&= \displaystyle{\sum_{\substack{\alpha\in v^t\Lambda_I^{\mathbf{e}_j}\\s_I(\alpha)\neq v^1,v^2}}} s_I(\alpha)(\N+\mathbf{e}_j)+\displaystyle{\sum_{\alpha\in v^t\Lambda_I^{\mathbf{e}_j}v^1}} s_I(\alpha)(\N+\mathbf{e}_j)+\displaystyle{\sum_{\alpha\in v^t\Lambda_I^{\mathbf{e}_j}v^2}} s_I(\alpha)(\N+\mathbf{e}_j)\\
&= \displaystyle{\sum_{\alpha\in v^t\Lambda_I^{\mathbf{e}_j}}} s_I(\alpha)(\N+\mathbf{e}_j)\\
&= v^t(\N)~~(\text{using relations of}~T_{\Lambda_I}).\qedhere
\end{align*}
}
\end{proof}

\end{document}
